% Appendix: the audit and its diagnosis in detail. Moved from Section~\ref{sec:quality}. Sources:
% evidence/kaggle-scale-run-2026-10-24.md (X0), evidence/x8-run.md (X8 repeats),
% evidence/x1-underretrieval-run1.md and evidence/x1-underretrieval-run2.md (X1; run 2 is the
% reference for every X1 number), investigations/x0-x1-pgvector-gap.md and
% investigations/filtered-recall-gaps.md (mechanisms).

\section{The audit and its diagnosis in detail (X0 and X1)}\label{app:quality}

This appendix holds the full audit, the full diagnosis of Qdrant and pgvector, and the design of the diagnosis grid,
which Section~\ref{sec:quality} summarizes. Appendix~\ref{app:ledger} lists the hypotheses.

\subsection{The audit}\label{app:audit}

\begin{table}[tbp]
\centering
\small
\begin{tabular}{@{}lrrrrrr@{}}
\toprule
 & \multicolumn{3}{c}{Four runs} & \multicolumn{2}{c}{Overlap@5, first run} & \\
\cmidrule(lr){2-4}\cmidrule(lr){5-6}
Store, under this layer's schema & Recall@5 & Recall@10 & Range@10 & judy & others & p50 (ms) \\
\midrule
FAISS (flat, exact) & 1.000 & 1.000 & 0.000 & 1.000 & 1.000 & 346--368 \\
Redis & 0.982 & 0.981 & 0.012 & 0.998 & 0.982 & 28--30 \\
Elasticsearch & 0.942 & 0.949 & 0.019 & 0.987 & 0.941 & 8--9 \\
OpenSearch & 0.942 & 0.935 & 0.004 & 0.998 & 0.941 & 7 \\
Weaviate & 0.899 & 0.902 & 0.021 & 1.000 & 0.876 & 7--8 \\
Chroma & 0.877 & 0.879 & 0.010 & 0.987 & 0.876 & 1432--1481 \\
LanceDB & 0.840 & 0.874 & 0.039 & 0.949 & 0.821 & 22--25 \\
pgvector, fast mode (2 runs) & 0.695 & 0.732 & 0.000 & 0.724 & 0.686 & 6.5--8.1 \\
pgvector, exact mode (2 runs) & 0.993 & 0.993 & 0.002 & & & 386--399 \\
Qdrant & 0.557 & 0.552 & 0.026 & 1.000 & 0.523 & 4.3--4.7 \\
\bottomrule
\end{tabular}
\caption{Retrieval quality in the enterprise scenario (Section~\ref{sec:e2e}), four runs of 1,105
questions (seeds 7, 11, 13, 17). Recall is the mean over runs; Range@10 is the largest minus the
smallest recall@10; p50 is the range of per-run medians. pgvector ran in one of two modes per run
(Section~\ref{sec:qual-pgvector}); each row averages the two runs in that mode. Overlap@5 is the
share of the exact top five that a store returned, by person type, from a re-analysis of the first
run (a fast-mode run for pgvector); ``judy'' are the least-privileged people (624 readable chunks).
pgvector with OPA follows pgvector's mode in each run. Settings: Appendix~\ref{app:settings}.}
\label{tab:qual-x0}
\end{table}

Table~\ref{tab:qual-x0} shows the spread that Section~\ref{sec:qual-audit} summarizes, at the settings of
Appendix~\ref{app:settings}. The rows compare configurations, not stores, and the corpus, in which only one of
twenty near-identical tenant copies of a document passes a person's filter, is a hard case for graph
indexes. Every store's loss is concentrated where the
readable set is far from the question. Of pgvector's first-run empty answers, eight were on red-team prompts and two on
questions about documents the person could not read. In the other two runs the
same configuration searched exactly, at fifty times the latency, and one of them still gave one
person one empty answer, as did another exact run of Section~\ref{sec:e2e}; we have not explained
these two. On targeted
questions in the first run, 59.6\% of Qdrant's top five were shared FiQA passages where exact search
returned none, and the model found the target document for only 39.9\% of targeted questions
(33.8\% to 39.9\% over the four runs).

\subsection{Qdrant}\label{app:qdrant}

Qdrant plans filtered search per segment from estimated filter cardinality~\cite{qdrant2026search}: a filter estimated below
a threshold (at the default 10,000\,KB and 768 dimensions about 3,333 points, by our arithmetic) runs as a plain, exact
scan over the payload index, and a broader one runs as filtered HNSW traversal, which in the source at v1.19.2 scores and
expands only neighbors that pass the filter. Extra graph edges that keep a
filtered subgraph connected are built per payload value, and only for values held by fewer points
than a cap of the point count divided by the average links per node, times four, about 19\% of the collection in the cited article's example~\cite{couzon2026filtered}. In the scenario every tenant
value (56\% of points) and
the employee and contractor groups (65\% and 61\%) exceed that.

\begin{table}[tbp]
\centering
\footnotesize
\begin{tabular}{@{}llrrr@{}}
\toprule
Layout & Search setting & Recall@10 & Range & p50 (ms) \\
\midrule
As run (one collection) & \texttt{hnsw\_ef} 128 (as run) & 0.656 & 0.086 & 7.1 \\
 & \texttt{hnsw\_ef} 64 & 0.600 & 0.074 & 6.0 \\
 & \texttt{hnsw\_ef} 512 & 0.758 & 0.120 & 11.3 \\
 & ACORN, max selectivity 0.4 or 1.0 & 0.808 & 0.127 & 7.3 \\
 & exact & 1.000 & n/a & 43.1 \\
Tenant field marked \texttt{is\_tenant} & as run & 0.651 & 0.025 & 6.9 \\
 & ACORN & 0.802 & 0.076 & 7.1 \\
Built as one segment & as run & 0.625 & 0.029 & 7.0 \\
 & ACORN & 0.777 & 0.027 & 7.2 \\
Per-tenant graphs (\texttt{is\_tenant}, $m = 0$) & as run & 0.952 & 0.001 & 6.6 \\
Shared passages in a separate collection & as run & 0.949 & 0.003 & 12.9 \\
\bottomrule
\end{tabular}
\caption{Qdrant 1.19.2 under the enterprise scenario's permission filters (X1 run 2, 1,489
questions, three index builds per layout, means over builds; Range is the range over builds; latency from
one CPU session). ACORN at the default gate (0.4) and at 1.0 gave identical results. Run 1 agreed within 0.05
(as run 0.619, ACORN 0.761, per-tenant graphs 0.943).}
\label{tab:qual-qdrant}
\end{table}

The pre-registered predictions for the hypothesis of Section~\ref{sec:qual-qdrant} all held, in both runs of the grid
(Table~\ref{tab:qual-qdrant} shows the second). They test recall by question category and by intervention; no connectivity or reachability was measured. In the first run, recall@10 for company documents among everyone but the least-privileged was 0.273, against 0.369 in the second. Per-tenant graphs were faster than the as-run layout (6.6 against 7.1\,ms). Raising \texttt{hnsw\_ef} from 128 to 512 stayed
within the band we had pre-registered as inconclusive and did not close the gap, as
a traversal trapped in the wrong region would predict. ACORN's 0.808 has a 95\% interval of 0.785 to 0.830 over people, while its range over the three builds is 0.127: read such intervals as conditional on the build.

One prediction failed. We expected ACORN's default gate, which engages when the estimated
selectivity is at most 0.4, to engage for some people and not others, because the planner's
independence estimate for a typical employee sits near 0.4. It engaged for everyone, and the gate
at 1.0 gave identical recall. So at \texttt{hnsw\_ef} 128 ACORN helps without reaching the layout change, and its gate does not misfire; with more effort it matches the change at lower throughput (Appendix~\ref{app:frontier}). A third hypothesis, that recall varies with the segment
layout, remains inconclusive: at the as-run setting recall@10 varied by up to 0.086 between three
builds of the same data (by up to 0.127 with ACORN),
and the layout built as one segment scored 0.030 lower (0.018 to 0.042), but in the end the as-run
collection also held its points in a single segment, so the difference is one between builds, not
layouts.

\subsection{pgvector}\label{app:pgvector}

pgvector 0.8 answers a filtered HNSW query by scanning the graph in order of distance and applying
the filter to each row the scan emits; with iterative scans enabled, it continues past
\texttt{hnsw.ef\_search} until enough rows pass the filter or a limit is
reached~\cite{pgvector2026pgvector}. The limit is \texttt{hnsw.max\_scan\_tuples} (20,000 by default)
or a memory bound of \texttt{hnsw.scan\_mem\_multiplier} times \texttt{work\_mem}. A person whose
readable set is small and far from the question can exhaust the budget before any readable row is
emitted, and the query then returns nothing, which is indistinguishable from having nothing to read.
That was our hypothesis for the empty answers of the audit. Testing it showed that a decision taken
before the scan matters more: whether PostgreSQL runs the HNSW scan at all.

\paragraph{The plan cache.}
The layer binds the filter's values and $k$ as parameters, so the search statement's text repeats
for every person with the same filter shape. psycopg prepares a statement on its sixth execution on
a connection~\cite{psycopg2026prepared}; PostgreSQL then plans five executions with their values and
afterwards uses a generic plan whenever that plan's estimated cost is not much higher than the
custom plans' average~\cite{postgresql2026prepare}, the trade-off studied as parametric query
optimization~\cite{ioannidis1997parametric}. A generic plan cannot see the \texttt{jsonb} values or
the \texttt{LIMIT}, assumes a selective filter, and reads the rows the GIN index returns and sorts
them: an exact search. The first diagnosis run recorded no executed plans, and most of its searches ran that one. The
second recorded the plan of every search (X1 run 2). Hypothesis P4 is confirmatory, pre-registered
before the first run; P5 is exploratory, specified in writing after the first run's results and
before the second ran.
Under the default configuration, after a vacuum, only 0.9\% of the searches on the scenario's questions used the HNSW
index, and recall@5 was 0.996 at a median of 362\,ms; planned with their own values, 98.1\% used
the HNSW index, with recall@5 of 0.680 at 22\,ms, against 0.691 in the audit's first run (P5
confirmed; P4, which had predicted that generic plans make no difference, falsified,
with recall@10 lower by 0.343 under custom plans).

One replay pair, on one build and in a fixed order, points to the state of the GIN index. Replaying the audit's pass right after loading,
when the index's pending list held 388 pages, PostgreSQL never switched: 98.1\% of searches used
HNSW, and recall@5 was 0.697 at 7.0\,ms, the audit reproduced. The same replay after
\texttt{VACUUM (ANALYZE)}, with the pending list empty, ran 2,148 generic executions, and recall@5 was 0.998 at
335\,ms. That command changes more than the pending list, so the pair does not isolate it. PostgreSQL's GIN cost estimate charges every pending page as a page read at the start of the
scan~\cite{postgresql2026selfuncs}, which is consistent with a generic GIN plan that looks expensive after a bulk load until a vacuum merges the list, and
with a generic plan that, once it has won, keeps winning on that connection. Autovacuum runs on its own
schedule, and the variation across the end-to-end runs would fit it, but we did not record
its timing there and cannot say what decided their modes: the same configuration ran fast by its median in three of eight runs, one of them a
mixture of fast and exact searches, and exact in five (Section~\ref{sec:e2e-safety}). A measurement
of pgvector through prepared statements can therefore
report one of two systems, 16 to 64 times apart in latency (16 on the grid, 23 to 64 across the eight end-to-end runs) and about 0.3 apart in recall, with no benchmark setting naming which.
The harness of Section~\ref{sec:qual-frontiers} plans
every search with its own values, and so can a deployment, with
\texttt{plan\_cache\_mode = force\_custom\_plan} or by not preparing the search statement.

\paragraph{The scan's budget.}
With the HNSW scan forced, the audit's loss came back on X1's full question set, which adds 24
people in a single group each: recall@10 of 0.656, and 8.9\% of retrievals with $k = 5$ empty, the
scan emitting a median of 20,181 rows on an empty answer. The budget hypothesis (P1) was
inconclusive as stated, and it needed a correction. Both limits bind: raising the tuple limit alone,
to 100,000 or a million, lowered the empty rate only to 7.3\%, because the 4\,MB memory bound then
stopped the scan; raising the memory alone changed nothing; raising both removed every empty
answer. The first run's report that the tuple limit alone sufficed came from searches that had in
fact run the exact plan. And 79.9\% of the empty answers came from questions whose first readable row
lay beyond half the scan's reach, short of the 95\% we had predicted.

Removing the empty answers did not remove the loss. With both limits raised, recall@10 was 0.780,
and 0.866 on a graph built with $m = 32$ and \texttt{ef\_construction} 200. The rest is
mis-retrieval: a relaxed-order scan stops at the first $k$ permitted rows it meets. The hypothesis
that this loss comes from shallow, relaxed scans (P2) was inconclusive. A deeper search
(\texttt{ef\_search} 1,000: $+0.069$), more candidates before the final sort (ten times $k$: $+0.095$)
and a better graph ($+0.066$) each helped, as predicted, and strict ordering hurt (0.497), but the
loss among questions with under a tenth of their 100 nearest chunks readable, 0.368, was not the
predicted twice the 0.188 elsewhere. The planner also
misjudged the people it served worst: for the least-privileged, its row estimate was a median 11.1
times their readable set (P3, inconclusive in the first run and confirmed in the second, the run that judges it), which would steer it to the HNSW scan where an exact plan over
a readable set of a few hundred chunks would return the right answer.

\subsection{The diagnosis grid}\label{app:x1grid}

\paragraph{pgvector.} The numbers in the paper are from run 2: three index builds with HNSW $m = 16$
and \texttt{ef\_construction} 64 and three with $m = 32$ and \texttt{ef\_construction} 200 (a core
grid), 1,489 questions (the scenario's 1,105 and 384 asked by 24 people in a single group each),
three timed passes each. Query-time settings: the scenario's settings with the planner's choice and
prepared statements (\emph{as run}); custom plans forced (\texttt{plan\_cache\_mode =
force\_custom\_plan}); no prepared statements; the HNSW scan forced (custom plans with sequential
and bitmap scans disabled), the reference for the scan-budget hypotheses, and every budget and depth
setting forced the same way: \texttt{hnsw.max\_scan\_tuples} of 20,000, 100,000 or 1,000,000 crossed
with \texttt{hnsw.scan\_mem\_multiplier} of 1, 8 or 16; \texttt{work\_mem} 64\,MB;
\texttt{hnsw.ef\_search} 40, 200, 400 and 1,000; an inner \texttt{LIMIT} of four or ten times $k$;
strict ordering; and exact search with index scans disabled. Two replays of the audit's pass ran
first on the first build, right after loading and after \texttt{VACUUM (ANALYZE)}; every other
setting ran after a vacuum. PostgreSQL's \texttt{auto\_explain} recorded the plan of every search of
the first pass. Run 1 built one pgvector index, skipped 12 settings for time, and recorded no
executed plans; it is reported in the released evidence.

\paragraph{Qdrant.} Five layouts (as run; tenant field marked \texttt{is\_tenant}; per-tenant graphs
with \texttt{payload\_m} 16 and $m = 0$; shared passages in a separate collection; one segment), three
builds each, with \texttt{hnsw\_ef} 64, 128 and 512, ACORN with maximum selectivity 0.4 and 1.0, and
exact search. Search paths per segment were read from Qdrant's telemetry.

\paragraph{Statistics.} 95\% intervals from 10,000 cluster-bootstrap resamples over people (seed
20261026); paired differences between settings on the same questions; verdicts computed by script
from the pre-registered criteria.
